\documentclass[11pt]{article}
\usepackage{fontspec}
\setmainfont{Times New Roman}
\setsansfont{Arial}
\setmonofont{Consolas}
\usepackage{amsmath,amssymb,mathtools}
\usepackage{geometry}
\usepackage{microtype}
\geometry{a4paper,margin=2.15cm}
\setlength{\parindent}{1.15em}
\setlength{\parskip}{0pt}
\emergencystretch=180em
\allowdisplaybreaks
\sloppy
\hbadness=10000
\vbadness=10000
\hfuzz=3pt
\vfuzz=10pt
\raggedbottom
\begin{document}% Unit: paper36_source_fidelity_v001.
% Language lane: Interslavic.
% Direct source: sources/paper36/source_fidelity/Noether_Paper36_ORIGINAL_JDMV39_GDZ600_R124plusP40_repaired_witness_v001.tex.
\section*{36. Diferenciranje idealov i diferenta}

\begin{center}
\emph{J. Ber. d. DMV39 (1930), s. 17}
\end{center}

E. Noether, Göttingen: Diferenciranje idealov i diferenta.

Pokazano jest, kako diferentu algebraičnogo čislovogo polja možno razuměti kako \emph{diferencialny kvocient} definujųcego \emph{ideala}. Že ta definicija sovpada s obyčnoj, slěduje iz strukturnyh teorem, interesnyh samih po sobě, ktore v analognom smyslu predstavjajų obobčenje Lagrangevoj interpolacijskoj formuly.

Podrobnějše izloženje ima se pojaviti v \emph{Mathematische Annalen}.
\end{document}